\section{Транспорт и форматы}\label{ACME.Transport}

\subsection{Пакеты HTTP}\label{ACME.Transport.HTTP}

Запросы на веб-узлы и ответы с них представляют собой пакеты
протокола HTTP. Пакеты должны защищаться при пересылке. Для организации защиты 
должен использоваться протокол, определенный в СТБ~34.101.65 (TLS версии 1.2) 
и СТБ~34.101.90 (TLS версии 1.3). В этом протоколе сервер ACME выступает в роли 
сервера TLS, клиент ACME~--- в роли клиента TLS.

% skip: ACME servers SHOULD follow the recommendations of [RFC7525] when
% configuring their TLS implementations.  ACME servers that support TLS
% 1.3 MAY allow clients to send early data (0-RTT).  This is safe
% because the ACME protocol itself includes anti-replay protections
% (see Section 6.5) in all cases where they are required.  For this
% reason, there are no restrictions on what ACME data can be carried in
% 0-RTT.

Роли сторон меняются и сервер ACME выступает в роли клиента при отправке 
вызова авторизации (см.~\ref{ACME.Chall}): клиента HTTP для вызова
\sstr{http-01} и клиента DNS для вызова \sstr{dns-01}.

% info: https://letsencrypt.org/docs/challenge-types/

Клиент ACME должен включать в свои HTTP-запросы заголовок \code{User-Agent}
(см.~\cite{HTTP}). В заголовке клиент дополнительно к названию и версии 
базового обеспечения HTTP должен указывать атрибуты программного обеспечения 
ACME.

Если сервер ACME разрешает доступ браузерных клиентов, то он должен поддерживать 
механизм CORS (Cross-Origin Resource Sharing, \cite{CORS}). Серверу в 
заголовке \code{Access-Control-Allow-Origin} своих HTTP-ответов следует 
указывать значение \sstr{*}.

Данные пакетов HTTP должны кодировать по правилам UTF-8,
определенным в~\cite{UTF8}. 
%
Если среди данных есть объект JSON, то двоичные данные этого объекта
должны кодироваться по правилам base64url~\cite{RFC4648} с исключением
завершающих символов \str{=}, а также пробелов, знаков табуляции и других 
дополнительных символов.
%
Объект JSON, в кодовом представлении двоичных полей которого которого 
встречается завершающий символ \str{=}, должен быть отклонен как 
некорректный.

\subsection{Ошибки}\label{ACME.Transport.Err}

% skip: Errors can be reported in ACME both at the HTTP layer and within
% challenge objects as defined in Section 8.  

При ошибке во время обработки запроса сервер ACME возвращает ответ с кодом 
статуса 4xx (ошибки клиента) или 5xx (ошибки сервера).

Если сервер сигнализирует об ошибке, то ему следует сообщить клиенту 
подробности проблемы в соответствии с~\cite{RFC7807}. Для этого сервер включает 
в ответ объект JSON, который содержит тип ошибки (поле \sstr{type}), ее краткое
описание (\sstr{title}), детальное описание (\sstr{detail}), другую информацию.
%
Типы ошибок в поле \sstr{title} должны начинаться с префикса 
\sstr{urn:ietf:params:acme:error:}. Далее в настоящем стандарте префикс 
опускается.

В таблице~\ref{Table.ACME.Err} определяются стандартные типы ошибок ACME, 
которые следует указывать в поле \sstr{type} подробностей проблемы.
%
Список типов ошибок в таблице не является исчерпывающим~--- сервер может 
указывать другие типы. 

% skip: Servers MUST NOT use the ACME URN namespace for errors not listed in 
% the appropriate IANA registry (see Section 9.6).  

Клиенту следует отображать поле \sstr{detail} подробностей проблемы.

\begin{table}[p]
\caption{Типы ошибок ACME}\label{Table.ACME.Err}
\begin{tabular}{|p{7cm}|p{9cm}|}
\hline
Тип & Описание\\
\hline
\hline
\sstr{accountDoesNotExist} &
Указанная в запросе учетная запись не существует
\\
\hline
%                                                                                                                           
\sstr{alreadyRevoked} &
Сертификат, указанный в запросе на отзыв, уже отозван
\\
\hline
%                                                                                                                           
\sstr{badCSR} &
Запрос на получение сертификата неприемлем 
\\
\hline
%                                                                                                                           
\sstr{badNonce} &
Клиент отправил некорректную синхропосылку
\\
\hline
%                                                                                                                           
\sstr{badPublicKey} &
Объект JWS был подписан открытым ключом, который сервер не поддерживает
\\
\hline
%                                                                                                                           
\sstr{badRevocationReason} &
Указанная причина отзыва сертификата не допускается сервером
\\
\hline
%                                                                                                                           
\sstr{badSignatureAlgorithm} &
Сервер не поддерживает алгоритм выработки ЭЦП, который использовался для 
формирования JWS
\\
\hline
%                                                                                                                           
\sstr{caa} &
Записи об авторизации УЦ запрещают выпускать сертификат
\\
\hline
%                                                                                                                           
\sstr{compound} &
Детали представлены в списке подробностей подпроблем
\\
\hline
%                                                                                                                           
\sstr{connection} &
Сервер не смог подключиться к объекту проверки
\\
\hline
%                                                                                                                           
\sstr{dns} &
Проблема с DNS-запросом во время проверки идентификатора
\\
\hline
%                                                                                                                           
\sstr{externalAccountRequired} &
В запросе должно быть указано значение поля \sstr{externalAccountBinding}
\\
\hline
%                                                                                                                           
\sstr{incorrectResponse} &
Полученный ответ не соответствует требованиям вызова
\\
\hline
%                                                                                                                           
\sstr{invalidContact} &
Контактный URL для учетной записи недействителен
\\
\hline
%                                                                                                                           
\sstr{malformed} & 
Запрос сформирован некорректно
\\
\hline
%                                                                                                                           
\sstr{orderNotReady} & 
Заявка не готова
\\
\hline
%                                                                                                                           
\sstr{rateLimited} &
Превышен лимит нагрузки сервера
\\
\hline
%                                                                                                                           
\sstr{rejectedIdentifier} & 
Сервер не будет выдавать сертификаты для заявленного идентификатора
\\
\hline
%                                                                                                                           
\sstr{serverInternal} & 
Внутренняя ошибка сервера
\\
\hline
%                                                                                                                           
\sstr{tls} & 
Ошибка TLS во время проверки
\\
\hline
%                                                                                                                           
\sstr{unauthorized} & 
У клиента нет достаточных полномочий
\\
\hline
%                                                                                                                           
\sstr{unsupportedContact} & 
В контактном URL учетной записи используется неподдерживаемая схема
\\
\hline
%                                                                                                                           
\sstr{unsupportedIdentifier} &
Проверяемый идентификатор не поддерживается
\\
\hline
%                                                                                                                           
\sstr{userActionRequired} &
Следует перейти по адресу, указанному в поле \sstr{instance}, и выполнить 
предлагаемые там действия
\\
\hline
\end{tabular}
\end{table}

Иногда серверу может потребоваться возвращать несколько ошибок в ответ на
запрос. Кроме того, серверу может потребоваться приписывать ошибки определенным
идентификаторам. Например, запрос на узел \code{newOrder} может содержать 
несколько идентификаторов, для которых УЦ не может выдать сертификат.

Объект JSON с подробностями проблемы может включать поле \sstr{subproblems}
со списком объектов JSON, которые описывают подробности подпроблем (частных 
проблем). Вложенные  объекты JSON имеют тот же формат, что и родительский объект. 
%
Вложенный объект может содержать поле \sstr{identifier}. Если поле присутствует, 
то в нем должен быть указан идентификатор ACME (см.~\ref{ACME.Res.Order}). 
При этом родительский объект не должен содержать поле \sstr{identifier}.
%
Типы ошибок в подробностях подпроблем не обязаны совпадать друг с другом и с 
типом ошибки в подробностях основной проблемы.

% skip: ACME clients may choose to use the "identifier" field of a subproblem 
% as a hint that an operation would succeed if that identifier were omitted.  
% For instance, if an order contains ten DNS identifiers, and the newOrder 
% request returns a problem document with two subproblems (referencing two of 
% those identifiers), the ACME client may choose to submit another order 
% containing only the eight identifiers not listed in the problem document.

\if0
\begin{codeblock}
HTTP/1.1 403 Forbidden
Content-Type: application/problem+json
Link: <https://example.com/acme/directory>;rel="index"

{
  "type": "urn:ietf:params:acme:error:malformed",
  "detail": "Some of the identifiers requested were rejected",
  "subproblems": [
    {
      "type": "urn:ietf:params:acme:error:malformed",
      "detail": "Invalid underscore in DNS name \"_example.org\"",
      "identifier": {
        "type": "dns",
        "value": "_example.org"
      }
    },
    {
      "type": "urn:ietf:params:acme:error:rejectedIdentifier",
      "detail": "This CA will not issue for \"example.net\"",
      "identifier": {
        "type": "dns",
        "value": "example.net"
      }
    }
  ]
}
\end{codeblock}
\fi

\subsection{Объект JWS}\label{ACME.Transport.JWS}

Запросы с непустым телом должны подписываться на личном ключе учетной записи
клиента. Подписанный запрос должен оформляться в виде объекта (контейнера)
JWS, определенного в~\cite{RFC7515} и описанного в СТБ~34.101.87 (Приложение Г).
%
Перед обработкой запроса сервер должен проверять подпись на открытом ключе 
учетной записи.

Объект JWS кроме базовых требований должен удовлетворять следующим 
дополнительным:

\begin{itemize}[label=--]
\item
должна использоваться плоская (flattened) JSON-сериализация объекта 
(см. \cite{RFC7515}). В частности, объект должен содержать только одну подпись 
(поле \sstr{signature});
\item
полезные данные (поле \sstr{payload}) не должны отсоединяться, их кодирование 
не должно отключаться (см.~\cite{RFC7797});
\item
защищенный заголовок (поле \sstr{protected}) должен содержать поля \sstr{alg}, 
\sstr{nonce}, \sstr{url} и одно из полей \sstr{jwk} или \sstr{kid}.
\end{itemize}

В поле \sstr{alg} заголовка указывается идентификатор алгоритмов ЭЦП,
которые будут использоваться для выработки и проверки подписей запросов.
%
Сервер ACME должен поддерживать идентификатор \sstr{BIGNS128}, 
установленный в СТБ~34.101.87 (приложение Г). Этот идентификатор описывает
алгоритмы \algname{bign-with-hbelt} (см.~\ref{CRYPTO.Sign}). Серверу ACME 
следует поддерживать идентификаторы \sstr{BIGNS192} и \sstr{BIGNS256}, которые  
описывают алгоритмы \algname{bign-with-bash384} и \algname{bign-with-bash512}.

Если сервер получил запрос с идентификатором, который не поддерживается,
то он должен дать ответ с кодом статуса 400 (некорректный запрос) и 
типом ошибки \sstr{badSignatureAlgorithm}.
%
Подробности проблемы в ответе должны содержать поле \sstr{algorithms} со 
списком поддерживаемых алгоритмов ЭЦП.

В поле \sstr{nonce} заголовка указывается синхропосылка. Ее клиент должен 
предварительно получить у сервера (см.~\ref{ACME.Transport.Reply}).
%
Синхропосылка должна быть представлена кодом base64url. При нарушении этого 
условия сервер должен отклонить JWS как некорректный. 
%
Если синхропосылка \sstr{nonce} имеет корректный формат, то сервер должен 
искать ее среди выпущенных.
%
Если синхропосылка не найдена, то сервер должен дать ответ с кодом статуса 400 
(недопустимый запрос) и типом ошибки \sstr{badNonce}.
%
Ответ должен включать заголовок \code{Replay-Nonce} с новой синхропосылкой, 
которую сервер примет при повторной отправке первоначального запроса (и, 
возможно, других запросов в соответствии с политикой сервера). Получив ответ, 
клиент должен повторить запрос, используя новую синхропосылку. 

В поле \sstr{url} заголовка повторяется адрес узла (URL), на который 
отправляется запрос. Таким образом блокируется подмена адреса.
%
Сервер должен сравнить адрес, на который получен запрос, со значением
в поле \sstr{url} и отклонить запрос в случае несовпадения адресов. 
%
В большинстве случаев адреса узлов выбираются сервером. Если это так, то 
клиент должен указывать в \sstr{url} посимвольно точный адрес, полученный от 
сервера. 
%
Серверу при обработке запроса на такой адрес следует проверять, что адрес точно 
указано в \sstr{url}. При обнаружении неточностей сервер должен отклонить запрос.

В поле \sstr{jwk} заголовка указывается открытый ключ учетной записи. Клиент 
должен использовать данное поле в своих запросах на узел \code{newAccount} 
(создание учетной записи) и может использовать в запросах на узел 
\code{revokeCert} (отзыв сертификата).

Открытый ключ в поле \sstr{jwk} представляет собой объект JWK, структура 
которого определена в~\cite{RFC7517}. 
%
Открытый ключ алгоритмов \algname{bign-with-hbelt}, \algname{bign-with-bash384} 
и \algname{bign-with-bash512} должен содержать три поля:
\begin{itemize}
\item
поле \sstr{kty} (строка)~--- тип ключа. Должно быть указано значение 
\sstr{OKP}, введенное в~\cite{RFC8037};
%
% info: https://datatracker.ietf.org/doc/html/rfc8037
%   A new key type (kty) value "OKP" (Octet Key Pair) is defined for public 
%   key algorithms that use octet strings as private and public keys. 

\item
поле \sstr{crv} (строка)~--- подтип ключа. Должен быть указан идентификатор 
алгоритмов ЭЦП: \sstr{BIGNS128}, \sstr{BIGNS192} или \sstr{BIGNS256};

\item
поле \sstr{x} (строка)~--- значение ключа. Кодовое представление 
по правилам base64url открытого ключа как строки октетов.
\end{itemize}

В поле \sstr{kid} заголовка указывается идентификатор ключей учетной записи.
Этот идентификатор назначает сервер и высылает клиенту в заголовке
\code{Location} своего ответа с узла \code{newAccount}. Клиент должен
использовать поле \sstr{kid} в запросах на узлы, отличные от \code{newAccount} 
и \code{revokeCert}. Клиент может использовать поле \sstr{kid} в запросах на 
узел \code{revokeCert}.

Сервер должен отклонить запрос, который содержит оба поля \sstr{jwk} и 
\sstr{kid}.

Если сервер поддерживает алгоритмы подписи, заданные в \sstr{alg}, но не 
поддерживает или отклоняет открытый ключ \sstr{jwk}, то он должен дать ответ с 
кодом статуса 400 (некорректный запрос) и типом ошибки \sstr{badPublicKey}. В 
подробностях проблемы должна быть указана причина отклонения открытого ключа.

Поскольку используется плоская JSON-сериализация объекта JWS, в заголовке 
\code{Content-Type} запроса как HTTP-пакета должно быть указано значение
\sstr{application/jose+json}. Если это не так, то сервер должен дать ответ с 
кодом статуса 415 (неподдерживаемый тип носителя).

\subsection{Запросы GET и POST-as-GET}\label{ACME.Transport.GET}

За исключением случаев, которые описываются ниже, сервер классифицирует 
запросы по HTTP-методам GET как ошибочные. В ответ на такой запрос сервер 
должен дать ответ с кодом статуса 405 (метод не разрешен) и типом ошибки 
\sstr{malformed}.

Для получения данных от сервера по схеме GET клиент должен использовать
HTTP-метод POST и включить в запрос объект JWS, в котором поле \sstr{payload}
представляет собой строку нулевой длины (\sstr{}). Такой запрос классифицируется
как POST-as-GET.

POST-запросы с пустым содержимым (и, следовательно, без объекта JWS) сервер 
должен обрабатывать так же, как GET-запросы.

Чтобы дать возможность клиентам подключаться к ACME, сервер должен разрешить 
GET-запросы к узлам \code{directory} и \code{newNonce} наряду с запросами 
POST-as-GET.

\subsection{Защита от повторов}\label{ACME.Transport.Reply}

Для защиты от атак на основе повторов сервер выпускает одноразовые синхропосылки
и рассылает их клиентам с требованием предъявлять в каждом подписанном запросе.
%
Клиент должен предъявлять синхропосылку в поле \sstr{nonce} заголовка объекта
JWS.
%
Сервер должен аннулировать синхропосылку сразу после предъявления клиентом, 
как будто бы синхропосылка ранее не выпускалась.

Сервер должен рассылать синхропосылки в заголовке \code{Replay-Nonce} 
своих ответов на POST-запросы клиентов. Синхропосылка должна включаться 
в каждый успешный ответ сервера и в каждый ответ об ошибке.
%
Заголовок \code{Replay-Nonce} не следует включать в запросы.

Сервер должен генерировать синхропосылки так, чтобы они повторялись лишь с 
незначительной вероятностью и были непредсказуемы для всех, кроме сервера.
%
Например, сервер может использовать в качестве синхропосылок случайные строки 
из 16 октетов.
%
Синхропосылка должна кодироваться по правилам base64url.
%
Клиент должен игнорировать некорректно закодированные синхропосылки 
в \code{Replay-Nonce}.

Настоящий стандарт не накладывает ограничений на то, как сервер определяет 
область действия синхропосылок. Клиенты могут считать, что имеется единый пул 
синхропосылок для всех запросов.
%
При этом в повторном запросе после ошибки \sstr{badNonce} 
(см.~\ref{ACME.Transport.JWS}) клиент должен указать синхропосылку из ответа    
об ошибке.
%
Сервер должны определять область действия синхропосылок достаточно широко, 
чтобы повторные запросы не требовались слишком часто.

\subsection{Контроль нагрузки}\label{ACME.Transport.Limit}

Сервер может контролировать нагрузку на свои узлы. При достижении предельного 
уровня нагрузки, то сервер в своих ответах должен указывать тип ошибки 
\sstr{rateLimited}.
%
Кроме этого, серверу следует включить в ответ заголовок \code{Retry-After} с 
информацией о том, когда текущий запрос может быть выполнен снова 
(см.~\cite{HTTP}).

% skip: If multiple rate limits are in place, that is the time where all rate 
% limits allow access again for the current request with exactly the same 
% parameters. 

% skip: In addition to the human-readable "detail" field of the error
% response, the server MAY send one or multiple link relations in the
% Link header field [RFC8288] pointing to documentation about the
% specific rate limit that was hit, using the "help" link relation
% type.

\subsection{Выпуск синхропосылок}\label{ACME.Transport.Nonce}

Для получения синхропосылки клиент отправляет HEAD-запрос на узел 
\code{newNonce} сервера (см.~\ref{ACME.Res.Dir}). 
%
Ответ сервера должен иметь код статуса 200 и включать заголовок 
\code{Replay-Nonce} с новой синхропосылкой. 

Кроме HEAD-запросов, сервер дополнительно должен обрабатывать GET-запросы на 
узел \code{newNonce}.
%
Ответ сервера на такой запрос должен иметь код статуса 204, включать заголовок  
\code{Replay-Nonce} с новой синхропосылкой и иметь пустое тело.

% use: https://errata.rfc-editor.org/eid6030/

\if0
HEAD /acme/new-nonce HTTP/1.1
Host: example.com

HTTP/1.1 200 OK
Replay-Nonce: oFvnlFP1wIhRlYS2jTaXbA

Cache-Control: no-store
Link: <https://example.com/acme/directory>;rel="index"
\fi

Чтобы предотвратить кэширование и повтор синхропосылок на стороне клиента,
сервер должен включать заголовок \code{Cache-Control} с директивой 
\sstr{no-store}.

